\section{InferCell}
\label{sec:design}

\IC{} is written in Julia \cite{bezanson2017julia}. It uses the SciML solvers for
ODEs and jump processes \cite{rackauckas2017differentialequations}. This section
gives the design ideas. The documentation gives the full interface
(see Software and data availability). \Cref{fig:architecture} shows the parts and how they
connect.

Each part answers one of the five obstacles of the introduction. Parameters that carry a
prior and a source give uncertainty a place in the model (\cref{sec:params}), and so does
a list of every reduction (\cref{sec:reductions}). Typed edges let ODE and jump modules
compose under one checked interface (\cref{sec:edges}), and the same gene expression runs
as an exact stochastic process or as mean-field ODEs (\cref{sec:twodesc}). The boundary
shares no parameter, so we say what inference across it can mean (\cref{sec:seam}). Two
inference routes then handle the latent paths and the cost: the exact route samples the
paths, and the cheaper mixed route integrates them out (\cref{sec:api}). Last,
\texttt{predict} turns a posterior into a prediction (\cref{sec:predictapi}).

\begin{figure*}[t]
\centering
\includegraphics[width=\linewidth]{fig_architecture.pdf}
\caption{\textbf{The \IC{} architecture.} Parameters carry a prior, a source and a label.
Modules declare their states, formalism and typed edges. The resolver checks the composition
and collects the declared reductions. The composed model runs as the exact hybrid or as one
mean-field ODE problem, and each description has an inference route. \texttt{predict} takes
a posterior from either route, a forward description and an optional intervention.}
\label{fig:architecture}
\end{figure*}

\subsection{Parameters are distributions with a source}
\label{sec:params}

Each parameter in \IC{} is an object that holds a value, a prior distribution, a
role and a record of its source. The record names the file, the table and the
identifier the value came from, and the values that other files give for the same
identifier. Each parameter also has one of four labels for how well its prior is
informed:
\begin{itemize}
  \item \emph{balanced}: the value and width come from a balanced parameter table;
  \item \emph{prior default}: the source gives no width, so the prior is wide;
  \item \emph{asserted}: the source gives no uncertainty, so we set the prior;
  \item \emph{not imported}: the parameter is not in the source.
\end{itemize}
The loader does not choose silently between two files that disagree. It stops and
asks for a decision, and the decision is then part of the record. A user can ask
the composed model which priors no measurement supports. That question has a
direct answer, because the answer is in the model and not in a supplement.

\subsection{Modules couple through typed edges}
\label{sec:edges}

A module declares its states, its parameters and its formalism: ODE or jump. It
also declares how it touches the rest of the cell. \IC{} has \pn{fw.edgekinds} kinds
of coupling edge (\cref{tab:edges}). Each kind says what crosses the edge and on
which clock. Mass and currency edges act inside the ODE solve. Catalytic edges
and deferred costs cross every second, at the handshake between the two blocks.
Rate-constant edges cross every minute, when the stochastic block rebuilds its rate
constants from the metabolite pools. The set of kinds is closed. A module cannot
invent a new kind of coupling without a change to the framework.

\begin{table}[t]
\caption{\textbf{The \pn{fw.edgekinds} edge kinds and their clocks.}}
\label{tab:edges}
\centering
\small
\begin{tabular}{@{}P{0.25\columnwidth}P{0.47\columnwidth}P{0.18\columnwidth}@{}}
\toprule
Edge & What crosses & Clock \\
\midrule
Mass & a shared continuous state & ODE solve \\
Currency & flux through a shared pool (ATP, GTP) & ODE solve \\
Catalytic & a protein count into a rate law & 1\,s \\
Deferred cost & jump-side cost, debited one step later & 1\,s \\
Rate constant & pools into stochastic rate constants & 60\,s \\
Volume & membrane proteins into cell geometry, and back & each step \\
Clamped & a dependence replaced by a constant & declared only \\
\bottomrule
\end{tabular}
\end{table}

Before anything runs, a resolver reads all the edges and builds the coupling graph.
It checks the composition for about \pn{fw.resolvererrors} kinds of error. Examples are a state that
two modules integrate, an edge that names a species no module owns, and a moiety
that is consumed but never produced. In complete mode it refuses to build a model
with an unowned state or a stranded moiety. A composition that runs is therefore
a composition that the framework has checked.

\subsection{Every reduction is declared}
\label{sec:reductions}

Any model of a cell leaves things out. In \IC{} each departure from the published
model is an object, not a comment. A reduction label has a category (one of
\pn{fw.redcats}, such as a clamp, a lumped step, a smoothed counter or an asserted
prior) and a reason. The framework collects labels from three places: the edges,
the parameter records and the modules. One call returns the full list for a
composed model. We report this list in \cref{sec:modelform}, and we measure the effect of only one entry on it.

\subsection{One interface, two descriptions}
\label{sec:twodesc}

Gene expression in \IC{} has one interface and two descriptions.

\paragraph{The exact hybrid.} The stochastic modules run as a jump process, and the
metabolic modules run as ODEs. A driver steps the two blocks in turn, as the
published model does. Every second it writes protein counts into the metabolic rate
laws, steps the ODEs, debits the costs of the stochastic events, and steps the jump
process. Every minute it rebuilds the stochastic rate constants from the pools. We deliberately did not wrap this in one solver object. A single jump problem over an ODE
problem would let the propensities follow the pools continuously. That is a
different model from the published one. The driver can also record a stochastic path
and replay it exactly. A replay matches the original run bit for bit at every
handshake. The exact sampler (\cref{sec:sampler}) depends on this.

\paragraph{The mean-field swap.} The gene-expression network of the model is first
order. So, given the rate constants, the equations for the means of the transcript and protein counts close exactly, as linear ODEs. \IC{} derives these equations from the same propensities
and update rules that the jump modules use, so the two descriptions cannot drift
apart. With the swap, the composed model becomes one ODE problem with callbacks for the
one-minute rebuild. Forward-mode automatic differentiation
\cite{revels2016forwardmode} passes through this problem and its callbacks. We test
the gradients against central differences to a relative error of $\pn{fw.adtol}$.

\paragraph{What is and is not differentiable.} The exact hybrid is not
differentiable, because the published model is not. The published model clips a
drained pool at zero and rounds pools to whole molecules. Both make the model a step
function of its parameters. The exact sampler therefore uses no gradients. The
mean-field description gives gradients for the parts of inference that need them.
The cost of the swap is a declared reduction, and we measure it
(\cref{sec:jensen}).

\subsection{Inference across a boundary with no shared parameter}
\label{sec:seam}

Three problems are often treated as one.
\begin{itemize}
  \item \emph{Local inference:} what do a module's own data say about its own
  parameters?
  \item \emph{Uncertainty propagation:} given the uncertain output of one module,
  how uncertain is the next module?
  \item \emph{Cross-module inference:} do the data of one module inform the
  parameters of another, when the two share no parameter?
\end{itemize}
The third problem is the hard one. Write $u$ for the trajectory at the interface: enzyme
counts and drains in one direction, pools in the other. For a pair of modules $A \to B$ with
data $D_A$ and $D_B$, the forward message is $p(u \mid D_A)$. The backward message is a
likelihood over the interface. For a one-way coupling, with $D_A$ and $D_B$ independent given $u$, \begin{equation}
  \lambda_B(u) = \int p(D_B \mid \theta_B, u)\, p(\theta_B)\, \mathrm{d}\theta_B .
\end{equation}
The joint posterior over $u$ is proportional to $p(u \mid D_A)\,\lambda_B(u)$. With feedback, as here, the messages alternate rather than close in one pass. A
\emph{statistical cut} drops $\lambda_B$ \cite{plummer2015cuts}. It changes the
target of inference, and data from $B$ can no longer inform $A$. A \emph{dynamical
reduction} instead removes a weak edge from the simulator and keeps full Bayesian
inference. For cross-module inference, only the second kind keeps the information we
want: metabolite data that inform a gene-expression parameter.

The standard modular methods pass messages through shared parameters: cuts and
semi-modular inference \cite{plummer2015cuts,carmona2020semimodular}, and Markov
melding \cite{goudie2019joining}. In the reduced model of \cref{sec:model} no parameter is shared, so
these methods have nothing to pass while we condition on the stochastic path. The
exact sampler does what is left: it conditions on the path, and it updates the path
with all the data. In this sense the exact sampler is modular inference with nothing
dropped. When the paths are integrated out, a gene-expression parameter appears in
both the transcript factor and the metabolite factor. The mixed route uses this
structure.

The framework makes the third question one that a modeller can ask of any composition. In
this paper we take only a first look at it, on M0 (\cref{sec:info}). Our claims rest on
calibration, and not on a measured transfer of information across the boundary.

\subsection{Inference attaches to each data stream}
\label{sec:api}

A data stream declares its observation model. In this paper there are two streams:
exact transcript counts per cell, and bulk metabolite measurements of the mean over
many cells. \IC{} has two inference routes.

\begin{itemize}
  \item \textbf{The exact route.} A three-block Gibbs sampler conditions on the
  stochastic path. It updates the paths by conditional sequential Monte Carlo
  \cite{andrieu2010particle}, the metabolic constants by slice sampling, and the
  gene-expression constants from exact one-dimensional conditionals
  (\cref{sec:sampler}).
  \item \textbf{The mixed route.} The paths are integrated out. The transcripts keep a birth--death likelihood given the mean-field pools, and the bulk metabolites get a Gaussian likelihood centred on the mean-field model, in the spirit of a synthetic likelihood \cite{wood2010statistical,price2018bayesian}. The route uses the gradients of the swap with the No-U-Turn
  Sampler \cite{hoffman2014nouturn}. It is an approximation, and
  \cref{sec:posteriors} compares it with the exact route on M0.
\end{itemize}

One function, \texttt{infer}, runs inference, and \texttt{inference\_route} chooses
the route from the composition and the data. An all-ODE composition with
time-series data takes the No-U-Turn Sampler, and an all-jump composition takes
approximate Bayesian computation. A hybrid composition with cell and bulk data has
the two routes above, and the user must name one, because one is a reference and the
other an approximation. Anything else is refused, with an error that says why. The
cell-data routes are built for M0's gene expression (\cref{sec:model}), and we have
demonstrated them only on M0.

\subsection{From posterior to prediction}
\label{sec:predictapi}

The function \texttt{predict} takes a posterior, a forward description and an optional
intervention, and it returns a predictive distribution for a chosen quantity. The composition
it predicts on may have parameters the fit never reached. We call those the held parameters
$\theta_H$, and the fitted ones the inferred parameters $\theta_I$. The held parameters get only
their priors and a viability rule $V$. The distribution is a cut in the sense of Plummer
\cite{plummer2015cuts}, with the first module a different model $M_{\mathrm{fit}}$ from the second,
\begin{equation}
  \begin{aligned}
  p_{\mathrm{cut}}(\theta_I, \theta_H \mid D, V)
  &= p(\theta_I \mid D, M_{\mathrm{fit}}) \\
  &\quad \times p(\theta_H \mid \theta_I, V),
  \end{aligned}
\end{equation}
and not a joint posterior under the second model. The data do not update
$\theta_H$, and $\theta_H$ does not feed back into $\theta_I$. A variance split then says how much
of the prediction's parameter uncertainty belongs to each block (\cref{sec:prediction}).

Which parameters are held is the user's choice. Freeing a fixed constant in its module's
constructor adds it to $\theta_H$, and \texttt{predict} draws it from its prior. Here a user
releases lactate dehydrogenase's forward constant and records GTP's time course:
{\footnotesize
\begin{verbatim}
ms = d11_models()
ms[1] = CentralGlycolysis(
  enzymes = :translated,
  free = [:kcatF_R_ENO, :kcatR_R_ENO,
          :kcatF_R_FBA, :kcatR_R_FBA,
          :kcatF_R_LDH_L, :kcatR_R_LDH_L])
p = predict(ms, posterior;
  horizon = COREA_CYCLE_S,
  completion = :prior, n_held = 20,
  n_cells = 4, trajectory = [:M_gtp_c])
\end{verbatim}
}
\noindent The first four names are the constants the reduced model already frees, and
\texttt{n\_held} is the number of draws of $\theta_H$ for each draw of $\theta_I$. Reverse
constants are derived from equilibrium constants and not drawn. \Cref{sec:release}
releases \pn{traj.nreleased} constants this way. A release is a sensitivity to priors and not a
posterior, since the fit held those constants at their nominal values.
